\section{Introduction}
\label{sec:intro}

Solar photovoltaic (PV) generation has grown rapidly across
Latin America, with Colombia reaching over 1~GW of installed capacity and
targeting 9~GW by 2030 under its National Energy Plan \cite{upme2023}.
Reliable short-term forecasts of global horizontal irradiance (GHI) are a
prerequisite for efficient operation of PV-integrated power systems: they
enable reserve scheduling, reduce imbalance penalties, and support real-time
dispatch and grid management. The growing share of variable renewable
generation makes this forecasting problem increasingly important: unlike
slowly-varying demand, cloud-induced irradiance variability is continuous,
stochastic, and correlated across spatially distributed generation assets,
and accurate, uncertainty-aware short-term forecasts provide operators with
the situational awareness needed to schedule reserves, trigger demand
response, and avoid costly imbalances.

In tropical countries such as Colombia, the forecasting challenge is compounded
by the high frequency and spatial heterogeneity of convective cloud systems,
which can reduce surface irradiance by more than 80\% within minutes and are
poorly captured by coarse-resolution numerical weather prediction (NWP) models.
Most published forecasting benchmarks, however, are built around temperate or
semi-arid climates, leaving tropical and high-altitude equatorial sites ---
where cloud dynamics differ substantially --- underrepresented in the
literature.

Satellite-based nowcasting methods bypass NWP limitations by inferring cloud
motion and optical depth directly from geostationary imagery. Instruments such as
the GOES-16 Advanced Baseline Imager (ABI) provide multi-spectral observations
every 10~minutes at sub-kilometre resolution, making them well-suited for
sub-hourly to hourly irradiance estimation and forecasting. The dominant paradigm
couples a spatial feature extractor (applied to a patch of satellite pixels
centred on the target site) with a temporal model that exploits the recent
history of observed GHI or satellite state \cite{zhang2015,reno2016}.

Deep learning has substantially advanced this paradigm. Convolutional neural
networks (CNNs) applied to satellite patches learn translation-invariant spatial
patterns associated with cloud regimes, and when combined with recurrent models
such as LSTMs they achieve state-of-the-art performance on benchmark datasets
\cite{alessandrini2015,huertas2018,paletta2021}. More recently, graph neural networks
(GNNs) have emerged as an alternative spatial encoder: by representing pixels as
nodes in a graph and learning neighbourhood aggregation rules, GNNs relax the
translation-invariance assumption of CNNs and can capture irregular spatial
dependencies introduced by non-uniform cloud structures \cite{hamilton2017,ji2022}.
However, systematic comparisons of GNN-based and CNN-based encoders for GHI
forecasting remain scarce, and none, to the best of our knowledge, address the
tropical Andean climate.

\subsection*{Contributions}

This paper makes the following contributions:

\begin{enumerate}
  \item We propose \textbf{GraphSAGE-LSTM}, a model that treats a GOES-16 satellite
        patch as a pixel graph, applies GraphSAGE neighbourhood aggregation
        \cite{hamilton2017} over it, and passes the resulting sequence of graph
        embeddings through an LSTM to produce a point forecast of GHI at horizons
        of 1, 3, and 6~hours. We compare a fixed unweighted nearest-neighbour
        graph against a learned, distance-weighted $k$-nearest-neighbour variant
        with a tunable neighbourhood size, isolating the contribution of
        \emph{learned spatial connectivity} to forecasting skill.

  \item We conduct a \textbf{rigorous empirical comparison} across two climatologically
        distinct Colombian sites, three forecast horizons, and multiple random seeds,
        including Bayesian hyperparameter optimisation via Optuna \cite{akiba2019}
        for GraphSAGE-LSTM, ResNet-LSTM, and a flat MLP spatial-ablation baseline,
        together with the graph-construction ablation described above.

  \item Preliminary single-seed results indicate that GraphSAGE-LSTM is
        \textbf{competitive with or superior to ResNet-LSTM} among the
        deep-learning models at multi-hour horizons at both sites, while the
        absolute skill of both models is markedly lower at the climatologically
        harder high-altitude site --- evidence that site-specific cloud
        dynamics, not just model architecture, govern attainable forecasting
        skill. The full multi-seed comparison, including the graph-construction
        ablation, is in progress (Section~\ref{sec:results}).

  \item We introduce a \textbf{disentangled uncertainty-quantification layer}
        that separately estimates the two sources of forecast uncertainty:
        an \emph{aleatoric} component, captured by a Gaussian
        negative-log-likelihood variance head that models irreducible cloud
        stochasticity, and an \emph{epistemic} component, captured by an
        ensemble of posterior weight samples obtained via Stochastic Gradient
        Langevin Dynamics (SGLD) \cite{welling2011}, which quantifies
        uncertainty due to limited data and model capacity
        \cite{kendall2017}. Unlike a single symmetric prediction interval,
        this decomposition tells an operator \emph{why} a forecast is
        uncertain, supporting differentiated, risk-stratified responses
        (e.g.\ wider reserve margins for high-epistemic-uncertainty regimes
        such as under-sampled cloud types, versus accepting irreducible
        variance during genuinely volatile but well-characterised convective
        afternoons).

  \item We establish a \textbf{reproducible GOES-16 benchmark for tropical
        Andean sites}, including data manifests, patch extraction, model code,
        Optuna studies, the graph-construction ablation, the SGLD/variance-head
        uncertainty-decomposition pipeline, and evaluation scripts, addressing
        the scarcity of published satellite-based GHI forecasting benchmarks
        --- with disentangled uncertainty quantification --- for Latin American
        and equatorial tropical contexts.
\end{enumerate}

\subsection*{Paper organisation}

Section~\ref{sec:related} reviews related work. Section~\ref{sec:methodology}
describes the datasets, model architectures, training protocol, and evaluation
metrics. Section~\ref{sec:results} presents quantitative results.
Section~\ref{sec:discussion} analyses key findings. Section~\ref{sec:conclusion}
concludes and outlines future work.
